\documentclass[11pt,a4paper]{article}
\usepackage[utf8]{inputenc}
\usepackage[french]{babel}
\usepackage[T1]{fontenc}
\usepackage{lmodern}
\usepackage{geometry}
\geometry{top=2.5cm,bottom=2.5cm,left=2.5cm,right=2.5cm}
\usepackage{xcolor}
\usepackage{hyperref}
\usepackage{titlesec}
\usepackage{enumitem}
\usepackage{booktabs}
\usepackage{array}
\usepackage{tcolorbox}
\usepackage{fancyhdr}

% Charte graphique corporate & moderne
\definecolor{primary}{RGB}{15, 23, 42}      % Navy / Slate 900
\definecolor{accent}{RGB}{37, 99, 235}      % Royal Blue
\definecolor{lightbg}{RGB}{248, 250, 252}   % Slate 50
\definecolor{darktext}{RGB}{51, 65, 85}     % Slate 700

\hypersetup{
    colorlinks=true,
    linkcolor=accent,
    urlcolor=accent,
    pdftitle={LuminaRH - Description du Projet},
    pdfauthor={Mamadou Ahmadou Ka}
}

% En-têtes et pieds de page
\pagestyle{fancy}
\fancyhf{}
\lhead{\small \textbf{LuminaRH} — Solution Numérique RH Industrielle}
\rhead{\small Hackathon 2026}
\cfoot{\thepage}

% Style des titres
\titleformat{\section}{\Large\bfseries\color{primary}}{\thesection.}{0.5em}{}[\titlerule]
\titleformat{\subsection}{\large\bfseries\color{accent}}{\thesubsection.}{0.5em}{}

\begin{document}

% Page de titre / En-tête principal
\begin{tcolorbox}[colback=primary,colframe=primary,arc=4mm,boxrule=0pt,top=15pt,bottom=15pt,left=15pt,right=15pt]
    \centering
    {\Huge \color{white} \textbf{LuminaRH}} \\[8pt]
    {\Large \color{white} \textbf{Plateforme digitale de gestion des Ressources Humaines pour l'industrie}} \\[12pt]
    {\small \color{white} Appel à Candidatures — Solutions facilitant des activités industrielles (ex: SOCOCIM)}
\end{tcolorbox}

\vspace{1em}

\begin{table}[h!]
\centering
\small
\begin{tabular}{ll}
\textbf{Porteur de projet :} & Mamadou Ahmadou Ka \\
\textbf{Établissement / Origine :} & Université de Thiès / Porteur de projet (StartUp) \\
\textbf{Email :} & \href{mailto:mahmadou.ka@univ-thies.sn}{mahmadou.ka@univ-thies.sn} \\
\textbf{Téléphone :} & +221 77 316 91 88 \\
\textbf{Thématique :} & Solutions facilitant des activités industrielles \\
\end{tabular}
\end{table}

\hrule
\vspace{1.5em}

\section{Résumé Exécutif}
\textbf{LuminaRH} est une plateforme numérique SaaS de gestion globale des Ressources Humaines adaptée aux réalités opérationnelles et réglementaires des entreprises industrielles du Sénégal (telles que SOCOCIM, ICS, ou le Port Autonome de Dakar). 

Grâce au pointage par reconnaissance faciale, au calcul automatisé de la paie selon le Code du travail sénégalais et la Convention Collective Interprofessionnelle, et à une gestion dématérialisée des congés et acomptes, LuminaRH transforme la gestion RH en un levier d'efficacité et de sérénité sociale pour l'industrie.

\section{Contexte et Problématique Industrielle}
Les grands sites industriels du Sénégal emploient de 500 à plus de 5 000 salariés\footnote{Source : Agence Nationale de la Statistique et de la Démographie (ANSD) — Recensement Général des Entreprises (RGE) \& Rapports annuels d'activités des grandes industries du Sénégal (SOCOCIM, ICS, PAD).} répartis en horaires décalés (3x8), avec un fort taux de travail posté et d'heures supplémentaires. Cependant, la gestion RH dans ces environnements fait face à des blocages sévères :

\begin{itemize}[leftmargin=1.5em]
    \item \textbf{Pointage manuel et fraude :} L'usage de pointeuses mécaniques à carte ou de feuilles de présence manuscrites engendre de la fraude au pointage (\textit{buddy punching}, estimé générer jusqu'à 5\% à 7\% de pertes sur la masse salariale brute selon les études du Bureau International du Travail — BIT\footnote{Organisation Internationale du Travail (OIT/BIT) — Études sur l’organisation du temps de travail et la productivité en Afrique subsaharienne.}) ainsi que des pertes de données liées aux environnements exigeants (poussière, chantiers).
    \item \textbf{Complexité de la paie locale :} Le calcul manuel sur tableurs Excel génère de fréquentes erreurs sur les cotisations sociales IPRES\footnote{Institution de Prévoyance Retraite du Sénégal (IPRES) — Barèmes des régimes Général (5,6\% salarial / 8,4\% patronal) et Cadre.} et CSS\footnote{Caisse de Sécurité Sociale du Sénégal (CSS) — Régimes des Prestations Familiales et Accidents du Travail.}, le calcul du retenu d'impôt sur le revenu (IR)\footnote{Code Général des Impôts du Sénégal (CGI) — Article 56 et suivant (Barème progressif avec quotient familial).}, le barème des heures supplémentaires majorées (+15\%, +40\%, +60\%)\footnote{Code du Travail du Sénégal (Loi n° 97-17) \& Convention Collective Nationale Interprofessionnelle (CCNI) — Article 27.} et l'indemnité légale de transport (fixée à 20 800 FCFA/mois)\footnote{Arrêté Ministériel n° 003926/MFPAI/DG-TSS relatif au montant de l'indemnité représentative de frais de transport au Sénégal.}.
    \item \textbf{Gestion des acomptes et prêts sociaux :} Absence d'automatisation dans le plafonnement légal des acomptes (plafonné à 50\% du salaire de base acquis selon l'Article L.113 du Code du Travail), entraînant un surcroît de travail administratif.
    \item \textbf{Coût et inadaptation des progiciels étrangers :} Les solutions internationales (SAP, Workday) sont excessivement chères et nécessitent de lourds développements spécifiques pour intégrer la législation fiscale et sociale sénégalaise.
\end{itemize}

\section{La Solution LuminaRH}
LuminaRH offre un écosystème web et mobile intégré pour automatiser l'intégralité du cycle de vie du collaborateur en milieu industriel.

\subsection{Modules Fonctionnels Clés}

\begin{itemize}[leftmargin=1.5em]
    \item \textbf{Pointage Biométrique \& Géolocalisé :} Système anti-fraude par reconnaissance faciale sur tablette/smartphone (avec cache local en cas de coupure réseau) et suivi des équipes en horaires 3x8.
    \item \textbf{Moteur de Paie Sénégalais Conforme :} Génération automatique des bulletins de paie intégrant :
    \begin{itemize}
        \item Cotisations IPRES (Régime Général \& Cadre) et CSS (Accidents du travail, Prestations familiales).
        \item Impôt sur le Revenu (IR) calculé selon les parts fiscales et le barème progressif national.
        \item Majoration automatique des heures de nuit, de dimanche et des jours fériés.
    \end{itemize}
    \item \textbf{Portail Self-Service Collaborateur :} Permet à chaque ouvrier ou cadre de consulter ses bulletins, faire des demandes de congés ou solliciter un acompte sur salaire en toute confidentialité.
    \item \textbf{Espace Administrateur \& Tableau de Bord :} Visualisation en temps réel de la masse salariale par unité de production, du taux d'absentéisme et gestion des départs (indemnités de licenciement, solde de tout compte, certificat de travail).
\end{itemize}

\section{Architecture Technique et Robustesse}
La plateforme repose sur une architecture moderne, modulaire, hautement sécurisée et conteneurisée :

\begin{table}[h!]
\centering
\small
\begin{tabular}{>{\bfseries}ll}
\toprule
Couche & Technologies Utilisées \\
\midrule
Frontend & React 19, TypeScript, Vite, TailwindCSS (Interfaces fluides \& Corporate) \\
Backend API & NestJS (Node.js), Architecture modulaire REST, OpenAPI/Swagger \\
Base de données & PostgreSQL 15, Prisma ORM, Séparation multi-schémas \\
Authentification & Keycloak SSO, Sécurité basée sur les rôles (RBAC) \\
Conteneurisation & Docker, Docker Compose, Redis (Gestion des sessions \& Cache) \\
\bottomrule
\end{tabular}
\end{table}

\section{Impact socio-économique et Valeur Ajoutée pour l'Industrie}
L'adoption de LuminaRH au sein d'une industrie comme la SOCOCIM permet :
\begin{enumerate}[leftmargin=1.5em]
    \item \textbf{Réduction de 90\% du temps de traitement de la paie} (de plusieurs jours à quelques minutes).
    \item \textbf{Élimination des erreurs financières} et des conflits sociaux liés aux bulletins inexacts.
    \item \textbf{Souveraineté des données :} Hébergement local des données RH des travailleurs sénégalais.
    \item \textbf{Accessibilité financière :} Modèle économique SaaS avantageux (abonnement par employé/mois) adapté aux PMI et grands groupes.
\end{enumerate}

\section{Références Légales \& Institutionnelles}
\begin{small}
\begin{enumerate}[leftmargin=1.5em]
    \item \textbf{ANSD (2023)} : \textit{Recensement Général des Entreprises (RGE)}, Agence Nationale de la Statistique et de la Démographie du Sénégal.
    \item \textbf{Code du Travail du Sénégal} : \textit{Loi n° 97-17 du 1er décembre 1997 portant Code du Travail de la République du Sénégal}.
    \item \textbf{CCNI du Sénégal} : \textit{Convention Collective Nationale Interprofessionnelle du Sénégal}.
    \item \textbf{Code Général des Impôts (CGI)} : \textit{Loi n° 2012-31 du 31 décembre 2012 portant Code Général des Impôts (Section IR/Retenue à la source)}.
    \item \textbf{Textes réglementaires IPRES \& CSS} : \textit{Guide des cotisations sociales employeurs du Sénégal (2024)}.
\end{enumerate}
\end{small}

\section{État d'Avancement}
Un MVP (Produit Minimum Viable) entièrement fonctionnel est d'ores et déjà développé et testé :
\begin{itemize}
    \item Base de données opérationnelle avec structure multi-schémas.
    \item Authentification centralisée avec gestion multi-rôles (Employé, Manager, RH/Admin).
    \item Interfaces d'administration et portails utilisateurs finalisés.
\end{itemize}

\end{document}
